\section{问题分析}

\subsection{问题一的分析}

问题一要求计算典型风光场景下园区运行的功率平衡和绿电指标。这是一个确定性的功率平衡核算问题——给定风光出力曲线和常规负荷曲线，在电解槽和合成氨装置满负荷连续运行的条件下，各时段的发电和用电功率均可精确确定。分析思路为：首先逐时计算风电、光伏出力与常规负荷、制氢氨负荷的差值，据此判定各时段需要网购电力还是有余电上网；然后汇总一日内的发电量、用电量、购电量、售电量，计算三个绿电直连指标和吨氨成本；最后将各指标与要求值对比，对不达标指标分析原因。本问题采用解析计算方法，无需建立优化模型。

\subsection{问题二的分析}

问题二要求园区制氨产能扩至72吨/日后，在电氢氨装置仅能全额开机或停机的条件下，寻找各产量水平(72-36吨/日)成本最低的生产时段安排。这是一个离散组合优化问题——每小时的生产决策为0-1变量（开机或停机），需要在满足总产量目标的前提下最小化日净电费（购电成本减售电收益）。分析思路为：首先注意到各时段的生产决策相互独立（无启停成本约束），因此可以分别计算每个时段开机的边际成本（相比不开机的费用变化）；然后对所有时段按边际成本从小到大排序，选择成本最低的$k$个时段开机，即可获得严格最优解。对于24种风光场景的统计分析，将每种场景作为独立算例重复求解，再按全满足、部分满足、全不满足三类对绿电指标达标情况进行分类统计。

\subsection{问题三的分析}

问题三在问题二的基础上，将电氢氨装置的运行方式从离散开关改为功率连续可调(下限10\%)。这是一个连续优化问题——各时段的制氢氨功率为连续变量，在满足总产量约束的前提下，通过优化功率在各时段的分配来最小化日总成本。分析思路为：该问题的目标函数是分段线性凸函数，可采用贪心分配策略求解——首先将功率优先分配给可再生盈余最大的时段（利用免费绿电，仅损失售电机会成本）；若仍有剩余产量需求，再按边际购电成本从小到大的顺序将功率分配到各时段。由于功率可连续调节，可以在各时段精确匹配可再生出力，预期将显著改善绿电指标。最后将计算结果与问题二进行对比，分析连续调节相比离散调节的优势。

\subsection{问题四的分析}

问题四要求分析园区离网运行时的能源自给状况并研究储能的最优配置。离网运行下园区与电网无电气连接，制氢氨功率必须跟随可再生出力变化，无法从电网购电补充。分析分三个层次：首先，在无储能条件下，计算各风光场景下制氢氨功率对可再生出力的跟随情况，统计产量、产能利用率和弃电率；其次，针对最大弃电场景，采用双层优化方法——外层枚举候选储能容量，内层在给定容量下优化逐时充放电调度，通过绘制Pareto曲线（吨氨成本与弃电率的权衡）寻找最优储能配置；最后，在相同产量需求下，对比离网模式与联网模式（调用问题三的连续调节模型）的经济性，包括吨氨成本、年总产量等指标。

\subsection{问题五的分析}

问题五是一个定性与定量相结合的政策分析问题。基于前四问的数值结果和电力系统运行的基本原理，从正反两方面分析高渗透率绿电园区对电网的影响：不利影响包括调峰压力、备用容量、电能质量等；有利影响包括促进消纳、降低碳排放、推动产业链等。政策建议基于定量论证（如本研究的成本数据和指标达标统计），按优先级排序给出，每条建议均有对应的数值依据支撑。
